\documentclass[11pt,a4paper]{article}
\usepackage[utf8]{inputenc}
\usepackage[T1]{fontenc}
\usepackage{amsmath,amssymb,amsthm,amsfonts}
\usepackage[hidelinks]{hyperref}
\usepackage{booktabs}
\usepackage{graphicx}
\usepackage{enumitem}
\usepackage[margin=1in]{geometry}

\newtheorem{definition}{Definition}
\newtheorem{conjecture}{Conjecture}

\title{RMD-Q: Restricted Module Decoding with Quadratic Constraint\\--- A New Post-Quantum Assumption with Lightweight KEM\\and Signature Constructions, Analysis, and Verified Prototype}
\author{nul0 (Future Proofs Tech) \\ \texttt{futureproofs@proton.me}}
\date{September 2026 \\ (DRAFT v0.1 --- research report, not a standard) \\
\small \textcopyright\ 2026 nul0 (Future Proofs Tech). Licensed under GPLv3 (see LICENSE).}

\begin{document}
\maketitle

\begin{abstract}
We introduce \emph{Restricted Module Decoding with Quadratic Constraint}
(RMD-Q), a new post-quantum hardness assumption that jointly couples a
noisy module-linear relation $b = As+e$ over $R_q = \mathbb{Z}_q[X]/(X^n+1)$
with a sparse multivariate-quadratic constraint $P(s)=0$ over the same
secret: breaking an instance requires a lattice decoder \emph{and} an MQ
solver together, neither alone suffices.  From one assumption we build both
a Fujisaki-Okamoto key-encapsulation mechanism and a Fiat-Shamir signature
with dual lattice/MQ openings, sharing ring arithmetic, XOF layer, and
codebase (crypto core $\approx 3.9$\,KB).  This report documents the
assumption (including a plantable key-generation distribution with measured
uniformity), exact decrypt-failure distributions by convolution (down to
$2^{-910}$ uncompressed, $2^{-235}$ at Kyber-style compression), dual-BKZ
and hybrid cost models validated within 3 bits on ML-KEM-512, a calibrated
kill-phase campaign (toy instances fall to brute force, MITM, LLL,
Gr\"obner, and msolve; larger instances stall or exhaust a 16-core/31\,GB
box), and a constant-time-shaped C implementation cross-checked
byte-for-byte against an independent Python reference (NTT backend with
exhaustively proven Barrett reduction, sanitizers and valgrind clean).
We make \emph{no} NIST category claims and \emph{no} deployment
recommendation: the assumption is conjectured, several analyses are
explicitly open, and Section~\ref{sec:nonclaims} lists everything this
report does \emph{not} claim.  All artifacts needed to reproduce every
number are described in Section~\ref{sec:repro}.
\end{abstract}

\noindent\textbf{Keywords:} post-quantum cryptography, module-LWE,
multivariate cryptography, planted problems, key encapsulation, digital
signatures, cryptanalysis, verified implementation

\tableofcontents
\newpage

%=====================================================================
\section{Introduction}
%=====================================================================

\subsection{Motivation}

The NIST post-quantum standards (ML-KEM~\cite{FIPS203}, ML-DSA~\cite{FIPS204},
SLH-DSA~\cite{FIPS205}) secure the general case, but leave two gaps this
work addresses as a \emph{research contribution, not a deployment proposal}:
(i)~ultra-constrained targets where every kilobyte of code and every
millijoule counts, and (ii)~assumption diversity beyond pure Module-LWE /
Module-SIS and hash-based designs (cf.\ the ongoing additional-signatures
process~\cite{IR8610}).

\subsection{Contributions}

\begin{enumerate}[leftmargin=*]
\item \textbf{A new assumption (Section~\ref{sec:problem}).}
  Search/Decision RMD-Q: sparse-short $s$ with $b=As+e$ \emph{and}
  $P(s)=0$ for a sparse homogeneous-quadratic system $P$.  We give a
  plantable key-generation distribution (always terminates, unlike
  rejection sampling at scale) with measured coefficient uniformity
  ($\chi^2=17.4$ vs.\ null mean $15$) and a documented leakage analysis.
\item \textbf{Two constructions from one assumption (Sections~\ref{sec:kem},
  \ref{sec:sig}).} An FO-KEM (full encaps/decaps with implicit rejection
  running in C at $n=256$) and a Fiat-Shamir signature with dual
  lattice/MQ openings (toy status, Section~\ref{sec:sig} states its limits).
  A found-by-testing soundness restriction (homogeneous-quadratics only)
  is documented with its root cause.
\item \textbf{Quantified analysis, no category claims
  (Section~\ref{sec:analysis}).} Exact failure distributions by convolution
  ($2^{-910}$ uncompressed; $2^{-235}$ at $(d_u,d_v)=(10,4)$); dual-BKZ
  model validated within 3 bits on ML-KEM-512 (20-bit optimism at 768,
  stated); guessing+lattice hybrid model (optimum: no guessing at these
  densities; wt lever quantified); joint XL model (MQ leg contributes
  negligibly in-model); Gr\"obner/msolve scaling curve with a 114\,GB OOM
  data point; 200\,000 C-speed failure trials, zero failures.
\item \textbf{Verified prototype (Section~\ref{sec:impl}).} Independent C
  and Python implementations cross-checked byte-for-byte on all KATs;
  NTT backend with exhaustively proven Barrett reduction; constant-time
  twins with automated timing smoke-test; ASan/UBSan/valgrind clean;
  Cortex-M4 cross-compile clean.
\end{enumerate}

\subsection{What this report is not (see Section~\ref{sec:nonclaims})}

Not a standard, not a submission, not deployment guidance, not a proof.
Security categories are never claimed; several analyses are explicitly
open; Section~\ref{sec:gaps} maps every claim to evidence and blockers.

%=====================================================================
\section{Preliminaries and notation}
\label{sec:prelim}
%=====================================================================

Let $R_q = \mathbb{Z}_q[X]/(X^n+1)$ with $n$ a power of two and $q$ prime;
our production ring is $(n,q)=(256,3329)$, the standard lattice ring shared
with ML-KEM~\cite{FIPS203} (deliberately: novelty is claimed only for the
joint assumption and constructions, never the arithmetic).
$S(\eta,w)$: each polynomial has at most $w$ nonzero coefficients, values
uniform over $[-\eta,\eta]$ then sparsified (reference
\texttt{sample\_sparse\_poly}); secrets only. $\mathrm{CBD}(\eta)$:
Kyber-style centered binomial, fixed-shape sampling; errors and encryption
randomness. $U(\mathbb{Z}_q)$: uniform; matrix $A$ expansion. Symmetric
layer is FIPS-202 SHAKE128/256 only~\cite{FIPS202}, with fixed domain
separation (matrix expand, FO coins/shared-secret, signature hashes,
implicit reject). KEM security target: IND-CCA2 via the modular
Fujisaki-Okamoto transform~\cite{HHK17}. Signatures follow Fiat-Shamir
with aborts~\cite{Lyu09,Lyu12}. BKZ cost models used:
$2^{0.292\beta}$ classical (BDGL sieve~\cite{BDGL16}), $2^{0.265\beta}$
quantum, $2^{0.2075\beta}$ paranoid.

%=====================================================================
\section{The RMD-Q assumption}
\label{sec:problem}
%=====================================================================

\begin{definition}[Search-RMD-Q]
Parameters $(n,q,k,\eta,w,\eta_e,t,m)$, $N=k{\cdot}n$.
\textsf{Gen} samples $A \gets U(R_q^{k\times k})$ (SHAKE128),
$s \gets S(\eta,w)^k$, $e \gets \mathrm{noise}(\eta_e)$,
$P \gets \mathrm{MQ}(t,m)$ conditioned on $P(s)=0$, and sets $b = As+e$.
Given $x=(A,b,P)$, find $s'$ with $\mathrm{wt}(s')\le k{\cdot}w$,
$\|s'\|_\infty \le \eta$, $\|As'-b\|_\infty \le B_e$, and $P(s')=0$
exactly. Decision-RMD-Q distinguishes $(A,As+e,P)$ from uniform.
\end{definition}

The assumption is that breaking an instance requires a lattice decoder
\emph{and} an MQ solver \emph{jointly}: lattice-only candidates fail
$P(s')=0$ with probability $\approx 1-q^{-t}$ each, while MQ-only roots
violate the noisy linear bound. No Shor structure exists (random module
matrix, random sparse MQ); Grover gives at most a square-root speedup over
the $C(N,w)\cdot(2\eta+1)^w$ sparse space.

\subsection{Plantable key generation}

Rejection-sampling $s$ until $P(s)=0$ succeeds with probability $\approx
q^{-t}$ per try -- impossible at scale. Our planted distribution samples
$s$ first, then per equation samples $m{-}1$ free terms (uniform supports
and coefficients over $[0,q)$) plus one fix monomial on
$\mathrm{support}(s)$ with solved coefficient
$c_{\mathrm{fix}} = -\mathrm{rest}\cdot \mathrm{ev}^{-1} \bmod q$, inserted
at a uniform slot; free coefficients are resampled (bounded) until
$c_{\mathrm{fix}}\ne 0$, with a pure-random escape hatch when all free
monomials vanish at $s$ (then uniform-random, no conditioning at all).
Given supports and $s$, coefficients are uniform over the hyperplane
$\{c : \langle c,\mathrm{evals}\rangle = 0\}$ minus a $1/q$ slice.
Measured: fix-coefficient $\chi^2 = 17.4$ on 209 samples over
$\mathbb{Z}_{17}^*$ (null mean 15); positions spread over $\mathrm{supp}(s)$.
Accepted residual leakage (pad linear terms on zero coordinates;
$1/q$ conditioning mass) is documented, not claimed away. First-ever joint
keygen at $(256, t{=}8)$ runs instantly with passing KEM round-trips.
This shifts random-MQ to \emph{planted}-MQ -- same family as planted
instances in MPCitH/code-based designs, but our sparse-planted distribution
needs dedicated analysis (Section~\ref{sec:analysis}).

%=====================================================================
\section{Constructions}
\label{sec:kem}
%=====================================================================

\subsection{KEM (FO-KEM over RMD-Q CPA-PKE)}
\label{sec:kemech}

The shape mirrors FIPS~203 (CPA-PKE + Fujisaki-Okamoto $\Rightarrow$
IND-CCA2~\cite{HHK17}) with the RMD-Q relation as trapdoor; $P$
participates in keygen and domain separation (ciphertexts carry no MQ
opening in v0.1 -- a documented limitation, not a security argument).
\textsf{KeyGen} expands $A$, samples sparse $s$, plants $P$, samples
CBD error $e$, sets $b=As+e$. \textsf{Encaps} derives coins as
$\mathrm{SHAKE256}(\text{coins}\|m\|\mathrm{pkhash})$, samples dense CBD
$(r,e_1,e_2)$, and outputs $u=A^Tr+e_1$,
$v=b^Tr+e_2+\mathrm{encode}(m)$ with $K$ derived alongside.
\textsf{Decaps} decodes $m'=v-s^Tu$, re-encrypts, byte-compares, and
returns $K$ or $\mathrm{SHAKE256}(\text{reject}\|\sigma\|ct)$ (implicit
rejection). Correctness needs per-coefficient noise $< q/4$; exact failure
distributions are computed in Section~\ref{sec:analysis}.

\subsection{Signature v1 (toy status)}
\label{sec:sig}

Same keypair shape; Fiat-Shamir with a \emph{dual} opening from one mask
$y$: $z_{\mathrm{lat}} = y + c_{\mathrm{poly}} s$ (sparse weight-$\tau$
challenge) for the lattice check $w' = Az_{\mathrm{lat}}-c_{\mathrm{poly}}b
\approx w$, and $z_{\mathrm{mq}} = y + c_{\mathrm{scalar}} s$ with bilinear
opening $h$ satisfying $P(z_{\mathrm{mq}}) = P(y) + c\,h$ for the MQ check.
\textbf{Load-bearing restriction (found by testing):} $P$ must be
homogeneous-quadratic -- with linear parts the opening gains a
$c^2P_{\mathrm{quad}}(s)$ defect and verification fails in general; toy
sparsity masked this (both parts vanished independently by luck). The
opening asserts $v_j\ne -1$. Unspecified in v0.1: prod challenge
$(\tau,\gamma_1,\beta)$ with size analysis, wire sizes, hedged policy,
rejection-rate analysis, n=256 C code, and any zero-knowledge property.

\subsection{Parameters}

\begin{center}
\begin{tabular}{lccccccccc}
\toprule
set & $n$ & $q$ & $k$ & $\eta$ & $w$/poly & $\eta_e$ & $t$ & $m$ & wire \\
\midrule
TOY-16 & 16 & 17 & 1 & 1 & 4 & 1 & 2 & 4 & raw \\
TOY-32 & 32 & 97 & 1 & 1 & 6 & 1 & 3 & 6 & raw \\
PROD-256 & 256 & 3329 & 2 & 2 & 128 & 2 & 8 & 12 & 12-bit \\
PROD-384 & 256 & 3329 & 3 & 2 & 160 & 2 & 10 & 12 & 12-bit (no KATs) \\
\bottomrule
\end{tabular}
\end{center}
The ring is shared with ML-KEM deliberately (NTT reuse, mature ring
cryptanalysis, comparability); $k{=}2/3$ matches ML-KEM-512/768 dimensions;
$w{=}128$ closes the modeled 5-bit gap to the Kyber reference at $\approx$
zero size cost (Section~\ref{sec:analysis}); $t/m$ supported by Gr\"obner
scaling and XL-infeasibility results below. Wire v0.1 is uncompressed
12-bit packing (pk $\approx$0.8\,KB, ct $\approx$0.8--1.1\,KB); compressed
$(10,5)$ FO path (800\,B ct, re-compress-then-compare decapsulation) proven
end-to-end in Python with a byte-exact C codec.

%=====================================================================
\section{Security analysis: models, measurements, kill-phase}
\label{sec:analysis}
%=====================================================================

All numbers below are directional (model tolerances stated); each result
points at its reproduction artifact (Section~\ref{sec:impl}).

\paragraph{Lattice leg.} Dual-BKZ model (Bai-Galbraith scaled dual;
$2^{0.292\beta}$/$2^{0.265\beta}$/$2^{0.2075\beta}$ costs) validates within
3 bits on ML-KEM-512 ($2^{115}$/$2^{104}$ vs.\ published $2^{118}$/$2^{107}$;
20-bit optimism at 768, stated). Ours: k=2 $\approx 2^{110}/2^{99}$,
$\approx$5--7 bits under same-$N$ reference in-model. fpylll cross-check of
the BKZ constant agrees within 0.3\% (conservative direction).

\paragraph{Hybrid guessing.} Position+value guessing + BKZ remainder:
optimum is $g{=}0$ at all densities (guessing uneconomical at 25\%
density); enumeration validated empirically; weight lever quantified
($w{=}256$ total closes the gap). Perfect-hints curve: $\approx$10 bits
lost per 32 leaked coordinates. Unhandled: hints-variants, BKW,
quantum walks.

\paragraph{Joint lattice+MQ.} XL with field equations counted
attacker-favorably is infeasible at prod scale (never closes at any degree
$\le 14$); lattice path binds in-model. Planted-structure-specific attacks
remain unmodeled -- the open cryptanalytic question.

\paragraph{Gr\"obner/msolve kill-phase.} Planted n=16: basis 0.3--2s;
full \emph{solve} $>$300\,s even at n=16 (bases cheap, solutions hard).
n=24 basis 180\,s, n=28 basis $>$1200\,s (bounded, unfinished), n=32 basis
exhausts 114\,GB on 16 cores (kernel-logged OOM -- memory wall, not timeout). Pure-MQ
(no bounds) terminates trivially everywhere but recovers nothing
(positive-dimensional variety). Planted vs.\ random show no measurable
non-genericity at n=16; Bardet bounds uninformative for this shape
(predicted 15/27 vs.\ observed 2--4). Toy instances additionally fall to
brute force, MITM (exact key, 23M checks in C), and LLL+Babai hybrids.

\paragraph{Failures.} Exact convolution distributions (FIPS-203 style):
$2^{-910}$--$2^{-512}$ per-KEM uncompressed across weights;
$2^{-268.6}$ at $(d_u,d_v)=(10,5)$, $2^{-235.4}$ at $(10,4)$ -- far under
the $2^{-138}$ bar. 200\,000 C-speed trials at n=256 (both paths): zero
failures (Clopper-Pearson 95\% upper $\approx 2^{-15}$).

%=====================================================================
\section{Implementation and verification}
\label{sec:impl}
%=====================================================================

Independent C and Python implementations cross-checked byte-for-byte on
all KATs (matrix expansion, sampling, ciphertexts, shared secrets,
implicit-reject keys, signature vectors). NTT backend for $(256,3329)$
(verified schedule: forward pairs equal direct CRT reduction; full
multiply equals schoolbook): $6.0\times$ multiply, $8.8\times$ matrix-vector
($\approx$25\,$\mu$s, full KEM $\approx$0.25\,ms host; $\approx$1.6M M4
cycles projected). Barrett reduction exhaustively proven over all 11M
products. Constant-time twins (masked arithmetic, Barrett-only, inverse
table) equivalence-tested with an automated objdump smoke check (zero
sign-jumps; reference still shows its jump, proving sensitivity); CT costs
$\approx$1.05--1.17$\times$. CBD sampling wired into FO paths; sparse
legacy retired behind \texttt{-DRMDQ\_NO\_LEGACY} (hardened-only build
fully green). ASan/UBSan/valgrind clean; strict warnings plus analyzer
silent; Cortex-M4 \texttt{-Werror} cross-compile clean (crypto core
$\approx$3.9\,KB; NTT scratch 24\,KB shared -- too fat for small MCUs,
streaming path retained). Sizes: pk $\approx$0.8\,KB, ct $\approx$0.8\,KB
uncompressed. One-command reproduction: \texttt{run\_all.sh}.

%=====================================================================
\section{Non-claims, gaps, and roadmap}
\label{sec:nonclaims}
\label{sec:gaps}
\label{sec:repro}
%=====================================================================

\textbf{Not claimed:} any NIST category; any reductionist proof (no ROM/QROM
sketches exist); certification of RMD-Q/planted-MQ; side-channel resistance
beyond measured timing smoke; any signature security level; deployment
readiness (no audit, single-threaded statics, variable-time Python
reference). Model numbers are directional with stated tolerances.

\textbf{Blocking gaps:} joint Gr\"obner degree-of-regularity at scale;
proper noisy-hint framework; BKW/quantum-walk analyses; fpylll full-attack
cross-check; exact failure confirmation via C trials at scale; compressed
wire KATs; signature prod parameters and proofs; formal object-level timing
proof; external review.

%=====================================================================
\section{Conclusion}
%=====================================================================

RMD-Q couples lattice and multivariate hardness in one assumption supporting
both KEM and signatures from a shared lightweight codebase, with
measurement-backed analysis and an explicit, checkable list of what remains
unknown. We invite attack on every open item -- that scrutiny is the point
of this report.


\bibliographystyle{plain}
\begin{thebibliography}{99}

\bibitem{FIPS203} National Institute of Standards and Technology.
  \emph{Module-Lattice-Based Key-Encapsulation Mechanism Standard},
  FIPS 203, 2024.

\bibitem{FIPS204} National Institute of Standards and Technology.
  \emph{Module-Lattice-Based Digital Signature Standard}, FIPS 204, 2024.

\bibitem{FIPS205} National Institute of Standards and Technology.
  \emph{Stateless Hash-Based Digital Signature Standard}, FIPS 205, 2024.

\bibitem{FIPS202} National Institute of Standards and Technology.
  \emph{SHA-3 Standard: Permutation-Based Hash and Extendable-Output
  Functions}, FIPS 202, 2015.

\bibitem{IR8610} G.~Alagic et al.
  \emph{Status Report on the Second Round of the Additional Digital
  Signature Schemes for the NIST PQC Standardization Process},
  NIST IR 8610, 2026.

\bibitem{HHK17} D.~Hofheinz, K.~H\"ovelmanns, E.~Kiltz.
  A modular analysis of the Fujisaki-Okamoto transformation.
  TCC 2017.

\bibitem{Lyu09} V.~Lyubashevsky.
  Fiat-Shamir with aborts: Applications to lattice and factoring-based
  signatures. ASIACRYPT 2009.

\bibitem{Lyu12} V.~Lyubashevsky.
  Lattice signatures without trapdoors. EUROCRYPT 2012.

\bibitem{BDGL16} A.~Becker, L.~Ducas, N.~Gama, T.~Laarhoven.
  New directions in nearest neighbor searching with applications to lattice
  sieving. SODA 2016.

\bibitem{CN11} Y.~Chen, P.~Q.~Nguyen.
  BKZ 2.0: Better lattice security estimates. ASIACRYPT 2011.

\bibitem{APS15} M.~R.~Albrecht, R.~Player, S.~Scott.
  On the concrete hardness of Learning with Errors.
  J.~Math.~Cryptology, 2015. (Estimator methodology.)

\bibitem{MAT22} M.~Matzov.
  Report on the Security of LWE: Improved Dual Lattice Attack, 2022.
  (Dual-hybrid attack framework; our model is a simplified instance.)

\bibitem{Bardet04} M.~Bardet, J.-C.~Faug\`ere, B.~Salvy.
  On the complexity of Gr\"obner basis computation of semi-regular
  overdetermined algebraic equations. ICPP 2004.

\bibitem{MSolve} J.~Berthomieu, C.~Eder, M.~Safey El Din.
  msolve: A library for polynomial system solving, 2021.
  \url{https://msolve.lip6.fr}

\bibitem{FPLLL} The FPLLL development team.
  FPLLL, a lattice reduction library (fpylll Python wrapper).
  \url{https://github.com/fplll/fpylll}

\bibitem{pqm4} M.~Kannwischer et al.
  PQM4: Post-quantum crypto library for the ARM Cortex-M4.
  \url{https://github.com/mupq/pqm4}

\bibitem{Kyber} R.~Avanzi et al.
  CRYSTALS-Kyber (version 3.02): Algorithm specification and supporting
  documentation, 2021. (NIST PQC Round 3.)

\bibitem{Dilithium} S.~Bai et al.
  CRYSTALS-Dilithium (version 3.1): Algorithm specification and supporting
  documentation, 2021. (NIST PQC Round 3.)

\end{thebibliography}

\end{document}
